\documentclass[12pt,a4paper]{article}
\usepackage[utf8]{inputenc}
\usepackage[serbian]{babel}
\usepackage{amsmath}
\usepackage{amsfonts}
\usepackage{graphicx}
\usepackage{booktabs}
\usepackage{hyperref}
\usepackage{geometry}
\geometry{top=25mm, bottom=25mm, left=30mm, right=25mm}

\hypersetup{
    colorlinks=true,
    linkcolor=black,
    filecolor=black,
    urlcolor=blue,
    citecolor=black,
}

\title{\bf Dark Patterns u dizajnu korisničkog interfejsa}
\author{Seminarski rad iz predmeta Računarstvo i društvo \\\\ [1cm] {\bf Student:} Aleksa Vukadinović \quad {\bf Indeks:} 103/2022}
\date{Maj, 2026. godine}

\begin{document}

\maketitle
\newpage

\tableofcontents
\newpage

\section{Uvod}
Savremeno društvo nalazi se u fazi duboke digitalne zavisnosti, gde se većina ekonomskih, socijalnih i administrativnih aktivnosti obavlja posredstvom veb-sajtova i mobilnih aplikacija. U idealizovanim narativima o razvoju računarstva, tehnologija se posmatra kao neutralan alat projektovan da optimizuje ljudsko delovanje, proširi slobodu izbora i demokratizuje pristup informacijama. Međutim, kako računarski sistemi postaju kompleksniji, a ekonomski modeli digitalne ekonomije agresivniji, interfejsi koje svakodnevno koristimo sve češće prestaju da budu neutralni posrednici. Umesto toga, oni postaju sofisticirani instrumenti za bihevioralni inženjering.

Jedan od najdominantnijih, a ujedno i najsuptilnijih oblika manipulacije u modernom digitalnom prostoru jesu ,,tamni obrasci'' (\emph{Dark Patterns}). Ovaj termin, koji je 2010. godine smislio britanski UI/UX dizajner Hari Branjul (\emph{Harry Brignull}), odnosi se na elemente korisničkog interfejsa koji su namerno i pažljivo projektovani sa ciljem da prevare ili navedu korisnike na donošenje odluka koje primarno idu u prilog komercijalnim interesima platforme, a često na direktnu štetu samog pojedinca \cite{mathur2019dark}.

Za razliku od tradicionalnih softverskih grešaka (tzv. \emph{bagova}) ili prosto lošeg dizajna koji nastaje usled nedostatka resursa ili nekompetentnosti programera, tamni obrasci su svesni i zahtevaju visok stepen programerske veštine, analitike podataka i duboko razumevanje ljudske psihologije. Oni iskoišćavaju ljudske kognitivne prečice, sistematske pristrasnosti i nedostatak pažnje kako bi kreirali asimetriju informacija u kojoj je korisnik sistemski hendikepiran.

Cilj ovog rada jeste analiza fenomena obmanjujućeg dizajna. Kroz analizu različitih istraživanja na ovu temu, kao i stvarnih primera iz života, rad će istražiti istorijske korene ovih praksi, njihovu tehnološku podlogu, pravne posledice (i propuste), kao i duboke socijalne posledice koje tamni obrasci ostavljaju na različite društvene grupe ljudi na osnovu njihove digitalne pismenosti.

\newpage
\begin{thebibliography}{99}

\bibitem{mathur2019dark}
Mathur, A., Acar, G., Friedman, M. J., Lucherini, E., Mayer, J., Chetty, M., \& Narayanan, A. (2019).
\emph{Dark Patterns at Scale: Findings from a Crawl of 11K Shopping Websites}.
Proceedings of the ACM on Human-Computer Interaction, 3(CSCW), 1-32.
\url{https://arxiv.org/abs/1907.07032}

\bibitem{brignull2010dark}
Brignull, H. (2010).
\emph{Dark Patterns: deceptions designed online}.
A List Apart, 318.
\url{https://alistapart.com/article/dark-patterns-deceptions-designed-online/}

\bibitem{gray2018the}
Gray, C. M., Kou, Y., Battles, B., Hoggatt, J., \& Toombs, A. L. (2018).
\emph{The Ethics of Undersigned Coercion: Types and Ethics of Dark Patterns}.
Proceedings of the 2018 CHI Conference on Human Factors in Computing Systems, 1-14.

\bibitem{zuboff2019age}
Zuboff, S. (2019).
\emph{The Age of Surveillance Capitalism: The Fight for a Human Future at the New Frontier of Power}.
PublicAffairs.

\end{thebibliography}

\end{document}
